% Source: https://artofproblemsolving.com/community/c3h2796294_2017_tx_mmc__texas_mathworks_math_contest
% Stable identity: aops-topic-2796294-q14
A point $(a,b)$ on the circle has tangent line $ax+by=5$. It passes through $(4,3)$ iff $4a+3b=5$.
\begin{asy}
usepackage("amsmath");
size(5cm);
defaultpen(fontsize(10pt));
import graph;
draw(circle((0,0),sqrt(5)));
pair P=(4,3),T1=(-0.4,2.2),T2=(2,-1);
draw(P--T1);draw(P--T2);draw((0,0)--T1,dashed);draw((0,0)--T2,dashed);
dot(P);dot(T1);dot(T2);dot((0,0));
label("$(4,3)$",P,NE);label("$\left(-\dfrac25,\dfrac{11}5\right)$",T1,NW);label("$(2,-1)$",T2,SE);label("$O$",(0,0),SW);
draw((-3,0)--(5,0),gray);draw((0,-3)--(0,4),gray);
\end{asy}

With $b=\dfrac{5-4a}3$: $9a^2+(5-4a)^2=45$, so $25a^2-40a-20=0$, $5a^2-8a-4=0$, $a=2$ or $a=-\dfrac25$. The points are $(2,-1)$ and $\left(-\dfrac25,\dfrac{11}5\right)$.

The tangent at $(a,b)$ has slope $-\dfrac ab$: at $(2,-1)$ the slope is $2$; at $\left(-\dfrac25,\dfrac{11}5\right)$ it is $\dfrac{2}{11}$. The smaller slope occurs at $(a,b)=\boxed{\left(-\dfrac25,\dfrac{11}5\right)}$.
